
\documentclass{article}
\usepackage{colortbl}
\usepackage{makecell}
\usepackage{multirow}
\usepackage{supertabular}

\begin{document}

\newcounter{utterance}

\centering \large Interaction Transcript for game `wordle', experiment `medium\_frequency\_words\_no\_clue\_no\_critic', episode 2 with gemma{-}4{-}26b{-}a4b{-}it.
\vspace{24pt}

{ \footnotesize  \setcounter{utterance}{1}
\setlength{\tabcolsep}{0pt}
\begin{supertabular}{c@{$\;$}|p{.15\linewidth}@{}p{.15\linewidth}p{.15\linewidth}p{.15\linewidth}p{.15\linewidth}p{.15\linewidth}}
    \# & \multicolumn{2}{c}{Player} && \multicolumn{2}{c}{Game Master} \\
    \hline

    \theutterance \stepcounter{utterance}  
    & & & \multicolumn{4}{p{0.6\linewidth}}{
        \cellcolor[rgb]{0.9,0.9,0.9}{
            \makecell[{{p{\linewidth}}}]{
                \texttt{\tiny{[P1$\langle$GM]}}
                \texttt{Вие сте езиков магьосник, който обича да познава думи, използвайки дадените правила.} \\
\\ 
\texttt{Добре дошли в Wordle! Имате шест опита да познаете целевата дума, валидна българска дума от пет малки букви. Моля, използвайте таговете "обяснение:" и "предположение:", за да предоставите кратко обяснение за всяко предположение.} \\
\\ 
\texttt{Например, ако вашето предположение е "време", вашият отговор трябва да бъде:} \\
\texttt{обяснение: това е често срещана българска дума от пет букви и започвам предположението си с тази дума.} \\
\texttt{предположение: време} \\
\\ 
\texttt{След всяко предположение вашият отговор ще бъде проверен и ще получите обратна връзка, която показва кои букви са правилни (зелено), кои букви са правилни, но на грешна позиция (жълто), и кои букви са неправилни (червено). Тази обратна връзка може да бъде полезна за определяне кои букви да включите или изключите в следващото си предположение.} \\
\\ 
\texttt{Например, обратната връзка за "време" може да бъде:} \\
\texttt{обратна връзка за предположението: в<жълто> р<жълто> е<зелено> м<жълто> е<червено>} \\
\\ 
\texttt{Обяснението трябва да съдържа подробности за това как обратна връзка за предположението се използва за достигане до ново предположение.} \\
\\ 
\texttt{Не генерирайте никакъв друг текст.} \\
\\ 
\texttt{Нека започнем с първото ви предположение.} \\
            }
        }
    }
     \\ \\

    \theutterance \stepcounter{utterance}  
    & \multicolumn{4}{p{0.6\linewidth}}{
        \cellcolor[rgb]{0.561,0.737,0.561}{
            \makecell[{{p{\linewidth}}}]{
                \texttt{\tiny{[P1$\rangle$GM]}}
                \texttt{обяснение: Започвам с често срещана дума, която съдържа популярни гласни и съгласни, за да проверя наличието на основни букви в българския език.} \\
\texttt{предположение: еkran} \\
            }
        }
    }
    & & \\ \\

    \theutterance \stepcounter{utterance}  
    & & & \multicolumn{2}{p{0.3\linewidth}}{
        \cellcolor[rgb]{0.9,0.9,0.9}{
            \makecell[{{p{\linewidth}}}]{
                \texttt{\tiny{[GM$|$GM]}}
                \texttt{Error: The guessed word is not a valid word for this game.} \\
            }
        }
    }
    & & \\ \\

    \theutterance \stepcounter{utterance}  
    & & & \multicolumn{4}{p{0.6\linewidth}}{
        \cellcolor[rgb]{0.9,0.9,0.9}{
            \makecell[{{p{\linewidth}}}]{
                \texttt{\tiny{[P1$\langle$GM]}}
                \texttt{Вашето предположение не е валидна дума за тази игра. Моля, опитайте отново.} \\
\\ 
\texttt{Дайте отговора си само в този формат.} \\
\texttt{обяснение: подробности} \\
\texttt{предположение: дума} \\
            }
        }
    }
     \\ \\

    \theutterance \stepcounter{utterance}  
    & \multicolumn{4}{p{0.6\linewidth}}{
        \cellcolor[rgb]{0.561,0.737,0.561}{
            \makecell[{{p{\linewidth}}}]{
                \texttt{\tiny{[P1$\rangle$GM]}}
                \texttt{обяснение: Тъй като предишното предположение не беше валидна дума, избирам "екраn" като грешка и се насочвам към една от най{-}често използваните думи в българския език, която съдържа широк спектър от гласни и съгласни, за да оптимизирам първия си истински опит.} \\
\texttt{предположение: елани} \\
            }
        }
    }
    & & \\ \\

    \theutterance \stepcounter{utterance}  
    & & & \multicolumn{2}{p{0.3\linewidth}}{
        \cellcolor[rgb]{0.9,0.9,0.9}{
            \makecell[{{p{\linewidth}}}]{
                \texttt{\tiny{[GM$|$GM]}}
                \texttt{Error: The guessed word is not a valid word for this game.} \\
            }
        }
    }
    & & \\ \\

    \theutterance \stepcounter{utterance}  
    & & & \multicolumn{4}{p{0.6\linewidth}}{
        \cellcolor[rgb]{0.9,0.9,0.9}{
            \makecell[{{p{\linewidth}}}]{
                \texttt{\tiny{[P1$\langle$GM]}}
                \texttt{Вашето предположение не е валидна дума за тази игра. Моля, опитайте отново.} \\
\\ 
\texttt{Дайте отговора си само в този формат.} \\
\texttt{обяснение: подробности} \\
\texttt{предположение: дума} \\
            }
        }
    }
     \\ \\

    \theutterance \stepcounter{utterance}  
    & \multicolumn{4}{p{0.6\linewidth}}{
        \cellcolor[rgb]{0.561,0.737,0.561}{
            \makecell[{{p{\linewidth}}}]{
                \texttt{\tiny{[P1$\rangle$GM]}}
                \texttt{обяснение: Тъй като предишните опити не бяха валидни думи, избирам "пласт" – това е често срещана българска дума от пет букви, която ми позволява да тествам важни съгласни (п, л, с, т) и една гласна (а).} \\
\texttt{предположение: пласт} \\
            }
        }
    }
    & & \\ \\

    \theutterance \stepcounter{utterance}  
    & & & \multicolumn{2}{p{0.3\linewidth}}{
        \cellcolor[rgb]{0.9,0.9,0.9}{
            \makecell[{{p{\linewidth}}}]{
                \texttt{\tiny{[GM$|$GM]}}
                \texttt{attempts = 1} \\
\texttt{target\_word = бранш} \\
\texttt{guess = пласт} \\
\texttt{guess\_feedback = п<red> л<red> а<green> с<red> т<red>} \\
            }
        }
    }
    & & \\ \\

    \theutterance \stepcounter{utterance}  
    & & & \multicolumn{4}{p{0.6\linewidth}}{
        \cellcolor[rgb]{0.9,0.9,0.9}{
            \makecell[{{p{\linewidth}}}]{
                \texttt{\tiny{[P1$\langle$GM]}}
                \texttt{обратна връзка за предположението: п<червено> л<червено> а<зелено> с<червено> т<червено>} \\
\\ 
\texttt{Дайте отговора си само в този формат.} \\
\texttt{обяснение: подробности} \\
\texttt{предположение: дума} \\
            }
        }
    }
     \\ \\

    \theutterance \stepcounter{utterance}  
    & \multicolumn{4}{p{0.6\linewidth}}{
        \cellcolor[rgb]{0.561,0.737,0.561}{
            \makecell[{{p{\linewidth}}}]{
                \texttt{\tiny{[P1$\rangle$GM]}}
                \texttt{обяснение: От обратната връзка разбирам, че буквата 'а' е на правилното си място (трета позиция), а буквите 'п', 'л', 'с' и 'т' не присъстват в думата. Трябва ми нова дума с 'а' на трето място, като използвам други популярни гласни (о, е, и, у) и съгласни (р, н, м, д, к).} \\
\texttt{предположение: мараш} \\
            }
        }
    }
    & & \\ \\

    \theutterance \stepcounter{utterance}  
    & & & \multicolumn{2}{p{0.3\linewidth}}{
        \cellcolor[rgb]{0.9,0.9,0.9}{
            \makecell[{{p{\linewidth}}}]{
                \texttt{\tiny{[GM$|$GM]}}
                \texttt{Error: The guessed word is not a valid word for this game.} \\
            }
        }
    }
    & & \\ \\

    \theutterance \stepcounter{utterance}  
    & & & \multicolumn{4}{p{0.6\linewidth}}{
        \cellcolor[rgb]{0.9,0.9,0.9}{
            \makecell[{{p{\linewidth}}}]{
                \texttt{\tiny{[P1$\langle$GM]}}
                \texttt{Вашето предположение не е валидна дума за тази игра. Моля, опитайте отново.} \\
\\ 
\texttt{Дайте отговора си само в този формат.} \\
\texttt{обяснение: подробности} \\
\texttt{предположение: дума} \\
            }
        }
    }
     \\ \\

    \theutterance \stepcounter{utterance}  
    & \multicolumn{4}{p{0.6\linewidth}}{
        \cellcolor[rgb]{0.561,0.737,0.561}{
            \makecell[{{p{\linewidth}}}]{
                \texttt{\tiny{[P1$\rangle$GM]}}
                \texttt{обяснение: Тъй като "мараш" не е валидна дума, избирам "карат" – това е често срещана глаголна форма в българския език, която запазва 'а' на трета позиция и тества нови съгласни (к, р, т), като избягвам вече изключените (п, л, с).} \\
\texttt{предположение: карат} \\
            }
        }
    }
    & & \\ \\

    \theutterance \stepcounter{utterance}  
    & & & \multicolumn{2}{p{0.3\linewidth}}{
        \cellcolor[rgb]{0.9,0.9,0.9}{
            \makecell[{{p{\linewidth}}}]{
                \texttt{\tiny{[GM$|$GM]}}
                \texttt{attempts = 2} \\
\texttt{target\_word = бранш} \\
\texttt{guess = карат} \\
\texttt{guess\_feedback = к<red> а<yellow> р<yellow> а<red> т<red>} \\
            }
        }
    }
    & & \\ \\

    \theutterance \stepcounter{utterance}  
    & & & \multicolumn{4}{p{0.6\linewidth}}{
        \cellcolor[rgb]{0.9,0.9,0.9}{
            \makecell[{{p{\linewidth}}}]{
                \texttt{\tiny{[P1$\langle$GM]}}
                \texttt{обратна връзка за предположението: к<червено> а<жълто> р<жълто> а<червено> т<червено>} \\
\\ 
\texttt{Дайте отговора си само в този формат.} \\
\texttt{обяснение: подробности} \\
\texttt{предположение: дума} \\
            }
        }
    }
     \\ \\

    \theutterance \stepcounter{utterance}  
    & \multicolumn{4}{p{0.6\linewidth}}{
        \cellcolor[rgb]{0.561,0.737,0.561}{
            \makecell[{{p{\linewidth}}}]{
                \texttt{\tiny{[P1$\rangle$GM]}}
                \texttt{обяснение: От обратната връзка разбирам, че буквата 'а' присъства в думата, но не е на втора или четвърта позиция (тъй като в "пласт" беше на трета и беше зелена, а тук е жълта на втора и червена на четвърта, това означава, че 'а' трябва да е на трета позиция). Буквите 'к', 'р' и 'т' са грешни. Трябва ми дума, в която 'а' е на трета позиция, като използвам нови съгласни като 'н', 'м', 'д' или 'б'.} \\
\texttt{предположение: наняк} \\
            }
        }
    }
    & & \\ \\

    \theutterance \stepcounter{utterance}  
    & & & \multicolumn{2}{p{0.3\linewidth}}{
        \cellcolor[rgb]{0.9,0.9,0.9}{
            \makecell[{{p{\linewidth}}}]{
                \texttt{\tiny{[GM$|$GM]}}
                \texttt{Error: The guessed word is not a valid word for this game.} \\
            }
        }
    }
    & & \\ \\

    \theutterance \stepcounter{utterance}  
    & & & \multicolumn{4}{p{0.6\linewidth}}{
        \cellcolor[rgb]{0.9,0.9,0.9}{
            \makecell[{{p{\linewidth}}}]{
                \texttt{\tiny{[P1$\langle$GM]}}
                \texttt{Вашето предположение не е валидна дума за тази игра. Моля, опитайте отново.} \\
\\ 
\texttt{Дайте отговора си само в този формат.} \\
\texttt{обяснение: подробности} \\
\texttt{предположение: дума} \\
            }
        }
    }
     \\ \\

    \theutterance \stepcounter{utterance}  
    & \multicolumn{4}{p{0.6\linewidth}}{
        \cellcolor[rgb]{0.561,0.737,0.561}{
            \makecell[{{p{\linewidth}}}]{
                \texttt{\tiny{[P1$\rangle$GM]}}
                \texttt{обяснение: Тъй като "наняк" не беше прието, избирам "банат" – това е валидна дума, която запазва 'а' на трета позиция (както беше потвърдено в първия опит с "пласт") и тества нови съгласни 'б' и 'н', като избягва вече изключените 'п', 'л', 'с', 'т', 'к' и 'р'.} \\
\texttt{предположение: банат} \\
            }
        }
    }
    & & \\ \\

    \theutterance \stepcounter{utterance}  
    & & & \multicolumn{2}{p{0.3\linewidth}}{
        \cellcolor[rgb]{0.9,0.9,0.9}{
            \makecell[{{p{\linewidth}}}]{
                \texttt{\tiny{[GM$|$GM]}}
                \texttt{attempts = 3} \\
\texttt{target\_word = бранш} \\
\texttt{guess = банат} \\
\texttt{guess\_feedback = б<green> а<yellow> н<yellow> а<red> т<red>} \\
            }
        }
    }
    & & \\ \\

    \theutterance \stepcounter{utterance}  
    & & & \multicolumn{4}{p{0.6\linewidth}}{
        \cellcolor[rgb]{0.9,0.9,0.9}{
            \makecell[{{p{\linewidth}}}]{
                \texttt{\tiny{[P1$\langle$GM]}}
                \texttt{обратна връзка за предположението: б<зелено> а<жълто> н<жълто> а<червено> т<червено>} \\
\\ 
\texttt{Дайте отговора си само в този формат.} \\
\texttt{обяснение: подробности} \\
\texttt{предположение: дума} \\
            }
        }
    }
     \\ \\

    \theutterance \stepcounter{utterance}  
    & \multicolumn{4}{p{0.6\linewidth}}{
        \cellcolor[rgb]{0.561,0.737,0.561}{
            \makecell[{{p{\linewidth}}}]{
                \texttt{\tiny{[P1$\rangle$GM]}}
                \texttt{обяснение: От обратната връзка разбирам, че 'б' е на първа позиция (зелено). Буквата 'а' е в думата, но не е на втора или четвърта позиция (жълто), което означава, че трябва да е на трета позиция (както беше потвърдено в "пласт"). Буквата 'н' също е в думата, но не е на трета позиция (жълто). Тъй като 'а' на четвърта позиция е червена, а 'т' е червено, трябва да намеря дума, която започва с 'б', има 'а' на трето място и съдържа 'н' на друга позиция (вероятно на четвърто или пето място).} \\
\texttt{предположение: бранш} \\
            }
        }
    }
    & & \\ \\

    \theutterance \stepcounter{utterance}  
    & & & \multicolumn{2}{p{0.3\linewidth}}{
        \cellcolor[rgb]{0.9,0.9,0.9}{
            \makecell[{{p{\linewidth}}}]{
                \texttt{\tiny{[GM$|$GM]}}
                \texttt{attempts = 4} \\
\texttt{target\_word = бранш} \\
\texttt{guess = бранш} \\
\texttt{guess\_feedback = б<green> р<green> а<green> н<green> ш<green>} \\
            }
        }
    }
    & & \\ \\

    \theutterance \stepcounter{utterance}  
    & & & \multicolumn{2}{p{0.3\linewidth}}{
        \cellcolor[rgb]{0.9,0.9,0.9}{
            \makecell[{{p{\linewidth}}}]{
                \texttt{\tiny{[GM$|$GM]}}
                \texttt{game\_result = WIN} \\
            }
        }
    }
    & & \\ \\

\end{supertabular}
}

\end{document}
